\subsection{Erstellung eines Kleinsignal-Ersatzschaltbilds}
% Alt: Aufgabe 20
\Aufgabe{
    \label{Aufg20}
    Zeichnen Sie das Kleinsignalersatzschaltbild der gegebenen Schaltung.
    \begin{figure}[H]
        \centering
        \input{Tikz/src/FigErsatzschaltbild1}\alttext{Die Abbildung zeigt eine Transistorschaltung. Ein Schaltbild zeigt eine Transistorschaltung mit einem PNP-Transistor. Der Emitter liegt unten an Masse, der Kollektor rechts führt zu der Ausgangsspannung \(U_\mathrm{A}\). Von der Versorgungsspannung \(U_\mathrm{V}\) oben führen zwei Widerstände nach unten: links ein \(75 \ k \Omega\)Widerstand zur Basis des Transistors, rechts ein \(500 \ \Omega\) Widerstand zum Kollektor. Links unten ist die Eingangsspannung \(U_\mathrm{E}\) mit Anschluss zur Basis.}
        \label{fig:FigErsatzschaltbild1}
    \end{figure}
}



\Loesung{
    \begin{figure}[H]
        \centering
        \input{Tikz/src/LsgErsatzschaltbild1}
        \alttext{Die Abbildung zeigt das Kleinsignalersatzschaltbild eines Bipolartransistors.
        Links ist die Eingangsspannung \(u_\mathrm{E}\) eingezeichnet, die zwischen dem Basisknoten „B“ und dem unteren Bezugsknoten gemessen wird. Zwischen Basis und Emitter befindet sich der Basis-Emitter-Widerstand \(r_{\mathrm{BE}}\). Zusätzlich ist zwischen Basis und Emitter der Widerstand \(R_2\) geschaltet.
        Zwischen Kollektor und Emitter ist der Kollektor-Emitter-Widerstand \(r_{\mathrm{CE}}\) eingezeichnet. Parallel dazu befindet sich eine gesteuerte Stromquelle. Deren Strom ist mit \(\mathrm{B} \cdot i_\mathrm{B}\) bezeichnet.
        Am Kollektorknoten „C“ ist der Lastwiderstand \(R_\mathrm{C}\) angeschlossen, der vom Kollektor zum Emitterknoten „E“ und weiter zur Masse führt. Die Ausgangsspannung \(u_\mathrm{A}\) wird zwischen Kollektor und Emitter gemessen.}
        \label{fig:LsgErsatzschaltbild1}
    \end{figure}
    Siehe: Abschnitt \ref{fig:ErsatzschaltbildEinesBipolartransistors} und \ref{fig:ErsatzschaltbilderEmitter-Kollektor-UndBasisschaltung}.
}